\section{Constructed Formal Vertices and Coefficientwise Correlators}
\label{sec:ch20-formal-vertices}

Section~\ref{sec:ch20-projected-products} treated density-wave and pairing
operators as \emph{literal products of finite-lattice fermion fields}.  We
could evaluate these products on the full CAR Fock space, and we saw why an
intermediate energy-budget projection introduces an explicit residual.  We
now ask a related but mathematically different question: can we construct a
model in which the oscillator expressions suggested by bosonization are
well-defined, possess an actual positive vacuum functional, and admit
exactly computable two-point functions?

A tempting shortcut would be to replace each fermion field by its
Mattis--Mandelstam expression, multiply exponentials, and invoke Gaussian
averaging.  This is not justified by the preceding finite-CAR calculation:
the projected product may have a nonzero residual, and compression does not
preserve every exponential identity.  Instead we shall \emph{construct a
second carrier}, with uncompressed canonical commutation relations (CCR),
independent charge shifts, and a formal-power-series parameter.  It makes
the algebraic steps behind the familiar vertex correlation formula exact,
while keeping the comparison with a finite interacting fermion state open.
We develop the uncompressed oscillator and charge construction here, giving
each definition and the arguments needed at the point where they first enter
the calculation.

\subsection{An uncompressed oscillator algebra and its positive vacuum}
\label{subsec:ch20b-oscillator}

Fix an integer \(h\geq 1\) and set \(L=2h\).  Spatial positions
\(x,y\) belong to \(\mathbb Z/L\mathbb Z\), and we use the canonical
Fourier character
\begin{equation}
 \zeta=\exp(2\pi i/L), \qquad
 \zeta^{m(x+L)}=\zeta^{mx}.
 \label{eq:ch20b-character}
\end{equation}
Let \(M\geq1\) be a \emph{finite} positive-mode cutoff and write
\(\mathcal I_M=\{(\nu,m):\nu\in\{R,L\},\ 1\leq m\leq M\}\).
The formal construction works for any finite \(M\); its intended
low-momentum comparison with the lattice usually has \(M\leq h\).
We use \(L\) for lattice length even when it appears as a subscript
for the left-moving species, where it is distinguishable by context.

Recall that an algebra is a vector space equipped with an associative
multiplication and a unit.  A \emph{star algebra} additionally has an
antilinear involution \(a\mapsto a^*\) satisfying
\((ab)^*=b^*a^*\).  The \emph{free} star algebra on symbols \(A_{\nu m}\)
and \(C_{\nu m}\) contains all finite linear combinations of their
noncommuting words.  Imposing an equation in this algebra means quotienting
by the two-sided, star-stable ideal generated by its difference: an
\emph{ideal} is a subspace closed under multiplication on either side by
arbitrary ambient words.  In the resulting quotient the imposed
relations hold for every word, not only when applied to a particular state.

\begin{definition}[The uncompressed oscillator algebra]
\label{def:ch20b-ccr}
Let \(\mathfrak A_{\mathrm{osc}}\) be the unital complex star-algebra
of finite words in \(A_{\nu m},C_{\nu m}\), modulo the relations
\begin{align}
 [A_{\nu m},C_{\eta n}]&=m\delta_{\nu\eta}\delta_{mn}I,
 & [A_{\nu m},A_{\eta n}]&=0,
 & [C_{\nu m},C_{\eta n}]&=0,
 \label{eq:ch20b-ccr}\\
 A_{\nu m}^*&=C_{\nu m},& C_{\nu m}^*&=A_{\nu m}.
 \nonumber
\end{align}
Here \([a,b]=ab-ba\).  There is no restriction on word degree or on
oscillator occupation number.
\end{definition}

The word \emph{uncompressed} is essential.  There are only \(2M\)
oscillator labels, but each creation generator may be applied any number
of times, so the representation space below is infinite-dimensional.
The trace of a commutator vanishes on a finite-dimensional space, whereas
the trace of \(mI\) does not: a faithful finite-dimensional realization of
\eqref{eq:ch20b-ccr} is therefore impossible.  This is the familiar
finite-dimensional obstruction from Chapter~\ref{ch:boson_fock}, not an
obstruction to the following polynomial representation.

\begin{proposition}[Concrete positive oscillator vacuum]
\label{prop:ch20b-vacuum}
There is a nonzero representation \(\pi_0\) of
\(\mathfrak A_{\mathrm{osc}}\) on the polynomial space
\(\mathcal P=\mathbb C[X_{\nu m}:(\nu,m)\in\mathcal I_M]\) and a
normalized, complex-linear, positive functional
\(\omega_0:\mathfrak A_{\mathrm{osc}}\to\mathbb C\), defined by
\begin{align}
 \pi_0(C_{\nu m})f&=X_{\nu m}f,
 &\pi_0(A_{\nu m})f&=m\,\frac{\partial f}{\partial X_{\nu m}},
 \label{eq:ch20b-polynomial-action}\\
 \langle X^\alpha,X^\beta\rangle_0
 &=\delta_{\alpha\beta}\prod_{(\nu,m)\in\mathcal I_M}
       \alpha_{\nu m}!\,m^{\alpha_{\nu m}},
 &\omega_0(a)&=\langle1,\pi_0(a)1\rangle_0.
 \label{eq:ch20b-gram}
\end{align}
The inner product is conjugate-linear in its first argument.  In
particular,
\(\omega_0(I)=1\) and \(\omega_0(a^*a)\geq0\) for every
\(a\in\mathfrak A_{\mathrm{osc}}\).
\end{proposition}
\begin{proof}
A multi-index \(\alpha\) has finitely many nonnegative integer
components.  Every weight in \eqref{eq:ch20b-gram} is strictly positive,
so the monomials form an orthogonal basis of a positive-definite
pre-Hilbert space; no completion is required.  Multiplication and
weighted differentiation satisfy
\([m\partial_{X_{\nu m}},X_{\eta n}]
 =m\delta_{\nu\eta}\delta_{mn}I\), and distinct multiplications or
differentiations commute.  Thus the defining commutator relations
annihilate every polynomial and the free-word action descends to the
quotient.  Moreover, if \(\beta=\alpha+e_{\nu m}\),
\begin{align*}
 \langle X_{\nu m}X^\alpha,X^\beta\rangle_0
 &=\prod_j\beta_j!\,m_j^{\beta_j}\\
 &=m(\alpha_{\nu m}+1)
         \prod_j\alpha_j!\,m_j^{\alpha_j}
 =\langle X^\alpha,m\partial_{X_{\nu m}}X^\beta\rangle_0.
\end{align*}
Both sides vanish if \(\beta\neq\alpha+e_{\nu m}\).  Hence creation
and annihilation are mutually adjoint on polynomials.  By reversing
words, \(\pi_0(a^*)\) is the algebraic adjoint of \(\pi_0(a)\).
Consequently
\[
 \omega_0(a^*a)=\langle\pi_0(a)1,\pi_0(a)1\rangle_0\geq0,
 \qquad\omega_0(I)=\langle1,1\rangle_0=1.
\]
In particular the quotient is not the zero algebra.  The vacuum
conditions \(\pi_0(A_{\nu m})1=0\) and
\(\langle1,\pi_0(C_{\nu m})f\rangle_0=0\) will determine its ordered
contractions.
\end{proof}

\begin{remark}[What positivity does and does not assert]
A normalized positive functional is \emph{linear}, not multiplicative.
Indeed a multiplicative scalar-valued map would send every commutator to
zero, contradicting \([A_{\nu m},C_{\nu m}]=mI\) for \(m\geq1\).
This polynomial vacuum is an algebraic state on an unbounded-operator
star algebra; we have not constructed a Weyl \(C^*\)-algebra or a
finite-dimensional fermionic density matrix.
\end{remark}

\subsection{Bogoliubov dressing and the charge-sector carrier}
\label{subsec:ch20b-dressed}

Let \(c,s\in\mathbb R\) satisfy \(c^2-s^2=1\).  When the remaining
coupling data matter, we denote their record by \(P\); in the present
section only \(c,s\) enter the oscillator formula.  Write
\begin{equation}
 g=(c-s)^2>0,\qquad g^{-1}=(c+s)^2,\qquad
 (c-s)(c+s)=1.
 \label{eq:ch20b-g}
\end{equation}
This Luttinger parameter \(g\) is not the Umklapp coupling \(g_U\) of
Chapter~21.  For each fixed positive mode, set
\begin{align}
 \widetilde C_{Rm}&=cC_{Rm}+sA_{Lm},
 &\widetilde C_{Lm}&=cC_{Lm}+sA_{Rm},
 \label{eq:ch20b-bog-creator}\\
 \widetilde A_{Rm}&=cA_{Rm}+sC_{Lm},
 &\widetilde A_{Lm}&=cA_{Lm}+sC_{Rm}.
 \label{eq:ch20b-bog-ann}
\end{align}
All four formulas retain the operator order inherited from the word
algebra.  The tilde refers to a change of generators, not to a
projection onto a smaller finite-energy space.

\begin{proposition}[Dressed CCR and positive state]
\label{prop:ch20b-dressing}
The substitution \(\beta_P(C_{\nu m})=\widetilde C_{\nu m}\),
\(\beta_P(A_{\nu m})=\widetilde A_{\nu m}\) extends to a star
algebra automorphism of \(\mathfrak A_{\mathrm{osc}}\).  Its inverse is
the same construction with \(s\) replaced by \(-s\).  The functional
\begin{equation}
 \omega_P=\omega_0\circ\beta_P^{-1}
 \label{eq:ch20b-omegaP}
\end{equation}
is normalized and positive.  Its dressed vacuum contractions are
\begin{equation}
 \omega_P(\widetilde A_{\nu m}\widetilde C_{\eta n})
 =m\delta_{\nu\eta}\delta_{mn},\quad
 \omega_P(\widetilde C_{\nu m}\widetilde A_{\eta n})=0,
 \quad\omega_P(\widetilde A\widetilde A)
 =\omega_P(\widetilde C\widetilde C)=0.
 \label{eq:ch20b-dressed-contractions}
\end{equation}
The last identities apply to every pair of indicated dressed modes.
\end{proposition}
\begin{proof}
As a representative calculation,
\[
 [\widetilde A_{Rm},\widetilde C_{Rn}]
 =c^2[A_{Rm},C_{Rn}]+s^2[C_{Lm},A_{Ln}]
 =(c^2-s^2)m\delta_{mn}I=m\delta_{mn}I.
\]
For the potentially nonzero cross-chirality annihilator commutator,
\[
 [\widetilde A_{Rm},\widetilde A_{Ln}]
 =cs\bigl([A_{Rm},C_{Rn}]+[C_{Lm},A_{Ln}]\bigr)=0.
\]
The other relations follow by exchanging \(R,L\), taking adjoints,
and using the bare relations.  The substitutions therefore respect the
relation ideal and star operation.  Applying the transformation with
\(-s\) composes the pairs
\((C_R,A_L)\) and \((C_L,A_R)\) by
\(\left(\begin{smallmatrix}c&-s\\-s&c\end{smallmatrix}\right)
  \left(\begin{smallmatrix}c&s\\s&c\end{smallmatrix}\right)=I\).
Thus \(\beta_P\) is invertible.  For every \(a\),
\[
 \omega_P(a^*a)
 =\omega_0\bigl((\beta_P^{-1}a)^*(\beta_P^{-1}a)\bigr)\geq0,
 \qquad\omega_P(I)=1.
\]
Since \(\beta_P^{-1}(\widetilde A)=A\) and
\(\beta_P^{-1}(\widetilde C)=C\), the contractions reduce to
\(\omega_0(AC)=m\) and
\(\omega_0(CA)=\omega_0(AA)=\omega_0(CC)=0\), mode by mode.
\end{proof}

The vertices must also change \emph{charge}.  For this reason the
oscillator algebra alone does not provide their carrier.  Let
\(\mathcal C=\mathbb C^{(\mathbb Z^2)}\) be the vector space of finite
linear combinations of kets \(\delta_N\), with
\(N=(N_R,N_L)\in\mathbb Z^2\) and
\(\langle\delta_N,\delta_{N'}\rangle=\delta_{NN'}\).
It is an algebraic direct sum, not the finite-band charge space of
Section~20.1.  The \(N_\nu\) act diagonally; hence all integer charge
labels are present.  Let \(e_R=(1,0)\) and \(e_L=(0,1)\).  Fix the
species ordering \(R<L\) and the reference half-band offset \(h\).
Define the signed bilateral shifts
\begin{align}
 F_\nu\delta_N&=P_\nu(N)\delta_{N-e_\nu},
 &P_\nu(N)&=(-1)^{\sum_{\eta<\nu}(h+N_\eta)+h+N_\nu-1},
 \label{eq:ch20b-signed-shift}\\
 F_\nu^*\delta_N&=P_\nu(N+e_\nu)\delta_{N+e_\nu},
 &Z_\nu(x)\delta_N&=\zeta^{N_\nu x}\delta_N.
 \label{eq:ch20b-phase-action}
\end{align}
The exponents of \((-1)\) are integers, even for negative charges.
Each \(F_\nu\) is unitary on this algebraic orthonormal basis:
its inverse and adjoint both have the displayed raising action.
Crucially, the left sign depends on the right charge, which encodes
the interspecies fermionic ordering.

\begin{theorem}[Anticommutation of Auxiliary Charge Shifts]
\label{thm:ch20b-shift-anticommute}
On the unbounded auxiliary module $\mathcal C = \mathbb C^{(\mathbb Z^2)}$, the charge shifts
satisfy exact mixed anticommutation and number shift relations:
\begin{equation}
 \{F_R, F_L\} = 0, \qquad
 \{F_R, F_L^*\} = 0, \qquad
 [N_\nu, F_\eta] = -\delta_{\nu\eta} F_\eta.
 \label{eq:ch20b-shift-car}
\end{equation}
\end{theorem}
\begin{proof}
Evaluate both operator orders on a basis vector $\delta_N$:
\begin{align*}
 F_R F_L \delta_N
  &= F_R \bigl(P_L(N) \delta_{N-e_L}\bigr)
   = P_L(N) P_R(N-e_L) \delta_{N-e_R-e_L}, \\
 F_L F_R \delta_N
  &= F_L \bigl(P_R(N) \delta_{N-e_R}\bigr)
   = P_R(N) P_L(N-e_R) \delta_{N-e_R-e_L}.
\end{align*}
By \eqref{eq:ch20b-signed-shift}, $P_R(N) = (-1)^{h+N_R-1}$, which is completely independent of $N_L$;
hence $P_R(N-e_L) = P_R(N)$.
On the other hand, $P_L(N) = (-1)^{(h+N_R) + (h+N_L-1)}$, which depends directly on $N_R$.
Shifting $N_R \to N_R - 1$ flips the parity:
\[
 P_L(N-e_R) = (-1)^{(h+N_R-1) + (h+N_L-1)} = - P_L(N).
\]
Therefore,
\[
 P_R(N) P_L(N-e_R) = - P_R(N) P_L(N) = - P_L(N) P_R(N-e_L),
\]
which proves $F_L F_R \delta_N = - F_R F_L \delta_N$ for all $N$, so $\{F_R, F_L\} = 0$.
The adjoint anticommutator $\{F_R, F_L^*\} = 0$ follows by an identical sign calculation.
Finally, since $F_\eta$ maps $\delta_N$ to $\delta_{N-e_\eta}$,
$[N_\nu, F_\eta]\delta_N = (N_\nu - \delta_{\nu\eta})F_\eta\delta_N - N_\nu F_\eta\delta_N = -\delta_{\nu\eta}F_\eta\delta_N$.
\end{proof}

\begin{theorem}[General Group 2-Cocycle Identity]
\label{thm:ch20b-cocycle-identity}
For arbitrary integer charge vectors $p, q \in \mathbb Z^d$, define the bilinear sign functional
\begin{equation}
 \varepsilon(p, q) := (-1)^{\sum_{1 \le i < j \le d} p_j q_i}.
 \label{eq:ch20b-cocycle-general}
\end{equation}
Then $\varepsilon$ satisfies the 2-cocycle condition:
\begin{equation}
 \varepsilon(p, q) \varepsilon(p+q, t) = \varepsilon(q, t) \varepsilon(p, q+t)
 \qquad (\forall p, q, t \in \mathbb Z^d).
 \label{eq:ch20b-cocycle-identity}
\end{equation}
\end{theorem}
\begin{proof}
Evaluating the exponent of the left-hand side modulo $2$ gives
$\sum_{i<j} [ p_j q_i + (p_j + q_j) t_i ] = \sum_{i<j} [ p_j q_i + p_j t_i + q_j t_i ]$.
The exponent of the right-hand side modulo $2$ is
$\sum_{i<j} [ q_j t_i + p_j (q_i + t_i) ] = \sum_{i<j} [ q_j t_i + p_j q_i + p_j t_i ]$.
These two sums are identical in $\mathbb Z$, proving the cocycle identity in $\mathbb C$.
\end{proof}

Let \(\mathfrak A_{\mathrm{ch}}\) be the \emph{generated} star algebra of
these operators and their finite sums and products, and set
\(\mathfrak A=\mathfrak A_{\mathrm{ch}}\otimes_{\mathbb C}
\mathfrak A_{\mathrm{osc}}\).  Operators from distinct tensor factors
commute.  On the algebraic tensor carrier
\(\mathcal C\otimes\mathcal P\), write
\begin{equation}
 \omega_{P,N}(a)=
 \big\langle\delta_N\otimes1,\,
      (\pi_{\mathrm{ch}}\otimes(\pi_0\circ\beta_P^{-1}))(a)
      (\delta_N\otimes1)\big\rangle.
 \label{eq:ch20b-tensor-state}
\end{equation}
A tensor product of these represented polynomial and finite-support
charge spaces has a positive-definite inner product.  The argument of
Proposition~\ref{prop:ch20b-vacuum} therefore gives
\(\omega_{P,N}(a^*a)\geq0\) and \(\omega_{P,N}(I)=1\).
The state depends on the \emph{source charge} \(N\).  It is not a
finite-CAR ground-state expectation: to make such a claim would require
a separately constructed comparison of states and operators.

\subsection{Hermitian phases and their commutators}
\label{subsec:ch20b-phases}

The opposite Fourier orientations of the two chiralities are important.
At position \(x\in\mathbb Z/L\mathbb Z\) define
\begin{align}
 Q_R(x)&=\sum_{m=1}^M\frac{1}{m}
 \bigl(\zeta^{-mx}\widetilde C_{Rm}
       +\zeta^{mx}\widetilde A_{Rm}\bigr),
 \label{eq:ch20b-QR}\\
 Q_L(x)&=\sum_{m=1}^M\frac{1}{m}
 \bigl(\zeta^{mx}\widetilde C_{Lm}
       +\zeta^{-mx}\widetilde A_{Lm}\bigr),
 \label{eq:ch20b-QL}\\
 X_{\mathrm{CDW}}(P,x)&=(c-s)\bigl(Q_R(x)+Q_L(x)\bigr),
 \label{eq:ch20b-Xcdw}\\
 X_{\mathrm{SC}}(P,x)&=(c+s)\bigl(Q_R(x)-Q_L(x)\bigr).
 \label{eq:ch20b-Xsc}
\end{align}
The factor \(1/m\) is fixed by our normalization
\([\widetilde A_{\nu m},\widetilde C_{\nu m}]=mI\).
The phases are oscillator operators; none of these symbols denotes a
finite-CAR field or its compressed image.

\begin{lemma}[Adjoints and cancellation of equal-channel commutators]
\label{lem:ch20b-hermitian-commuting}
For every \(x,y\),
\begin{equation}
 Q_\nu(x)^*=Q_\nu(x),\qquad
 X_J(P,x)^*=X_J(P,x),\qquad
 [X_J(P,x),X_J(P,y)]=0
 \quad(J\in\{\mathrm{CDW},\mathrm{SC}\}).
 \label{eq:ch20b-xcommute}
\end{equation}
\end{lemma}
\begin{proof}
Since \(\zeta\) has unit modulus,
\((\zeta^{-mx}\widetilde C_{Rm})^*
=\zeta^{mx}\widetilde A_{Rm}\); the other term is its adjoint.
The same applies to the left chirality.  Because \(c\pm s\) are real,
all the \(Q_\nu\) and \(X_J\) are Hermitian.

The CCR reduce the right-chiral commutator to
\begin{align*}
 [Q_R(x),Q_R(y)]
 &=\sum_{m=1}^M\frac1m
   \bigl(\zeta^{m(x-y)}-\zeta^{-m(x-y)}\bigr)I\\
 &=2i\sum_{m=1}^M\frac{\sin(2\pi m(x-y)/L)}m\,I.
\end{align*}
The reversed spatial orientation gives exactly the negative expression
for \([Q_L(x),Q_L(y)]\).  Mixed commutators vanish because distinct
dressed chiralities obey independent CCR.  Expanding a commutator of
either \(Q_R+Q_L\) or \(Q_R-Q_L\) with itself at two positions leaves
the \emph{sum} of the two same-chirality commutators, so the signs cancel.
Multiplication by \((c-s)^2\) or \((c+s)^2\) proves
\eqref{eq:ch20b-xcommute}.  Notice that a single \(Q_R\) does
\emph{not} generally commute with itself at different positions.
\end{proof}

We shall need a finite kernel rather than a limiting logarithm:
\begin{definition}[Finite logarithmic kernel]
\label{def:ch20b-D1}
\begin{equation}
 D_1(x,y)=\sum_{m=1}^{M}\frac{1-\cos(2\pi m(x-y)/L)}{m}.
 \label{eq:ch20b-D1}
\end{equation}
\end{definition}
Every sum here is finite.  Thus \(D_1(x,x)=0\),
\(D_1(x,y)=D_1(y,x)\), and \(D_1(x,y)\geq0\).
Since \(M\geq1\), the \(m=1\) term implies strict positivity if
\(x\not\equiv y\pmod L\).  No continuum limit is implicit in the
term ``logarithmic.''

\begin{proposition}[Exact difference variances]
\label{prop:ch20b-variance}
Put \(\Delta X_J=X_J(P,y)-X_J(P,x)\).  In the positive state
\(\omega_{P,N}\), which factorizes on oscillator and charge words,
\begin{equation}
 \omega_{P,N}(\Delta X_{\mathrm{CDW}})=0,
 \qquad\omega_{P,N}(\Delta X_{\mathrm{CDW}}^2)=4gD_1(x,y),
 \label{eq:ch20b-variance-cdw}
\end{equation}
\begin{equation}
 \omega_{P,N}(\Delta X_{\mathrm{SC}})=0,
 \qquad\omega_{P,N}(\Delta X_{\mathrm{SC}}^2)
       =4g^{-1}D_1(x,y).
 \label{eq:ch20b-variance-sc}
\end{equation}
\end{proposition}
\begin{proof}
Each \(Q_\nu\) is a finite linear combination of single dressed
creation and annihilation operators, so its expectation is zero by
\eqref{eq:ch20b-dressed-contractions}.  For the right chirality the
nonzero ordered contraction is
\[
 \omega_P(Q_R(x)Q_R(y))
 =\sum_{m,n=1}^M\frac{\zeta^{mx}\zeta^{-ny}}{mn}
   \omega_P(\widetilde A_{Rm}\widetilde C_{Rn})
 =\sum_{m=1}^M\frac{\zeta^{m(x-y)}}{m}.
\]
For the left chirality it is the conjugate expression,
\(\sum_m\zeta^{-m(x-y)}/m\).  In either chirality
\begin{align*}
 \omega_P\bigl((Q_\nu(y)-Q_\nu(x))^2\bigr)
 &=\omega_P(Q_\nu(y)^2)+\omega_P(Q_\nu(x)^2)\\
 &\quad-\omega_P(Q_\nu(y)Q_\nu(x))
       -\omega_P(Q_\nu(x)Q_\nu(y))\\
 &=2\sum_{m=1}^M\frac{1-\cos(2\pi m(x-y)/L)}m
 =2D_1(x,y).
\end{align*}
Mixed \(R,L\) two-point contractions vanish by the dressed vacuum
formula.  The variance of either \(\Delta Q_R+\Delta Q_L\) or
\(\Delta Q_R-\Delta Q_L\) is therefore \(4D_1\).
Multiply respectively by \((c-s)^2=g\) and
\((c+s)^2=g^{-1}\).  The charge tensor factor acts as the identity
throughout this calculation.
\end{proof}

\subsection{Charge words, formal series, and Gaussian coefficients}
\label{subsec:ch20b-formal}

It is instructive to separate the pieces of a vertex which change
charge from those which create oscillator excitations.  In the charge
algebra introduce the \emph{ordered} Klein--phase words
\begin{align}
 K_{\mathrm{CDW}}(x)
 &=F_R^*F_L\,\zeta^{(N_L-N_R-1)x},
 \label{eq:ch20b-Kcdw}\\
 K_{\mathrm{SC}}(x)
 &=F_RF_L\,\zeta^{(N_R+N_L)x}.
 \label{eq:ch20b-Ksc}
\end{align}
The rightmost diagonal phase acts on the source charge before either
shift.  Replacing the product by one with the phase on the left changes
its coefficient: these factors need not commute.  The star in
\(F_R^*\) is the adjoint shift, not complex conjugation of a scalar.

\begin{lemma}[Klein products at fixed source charge]
\label{lem:ch20b-klein}
For every \(N=(N_R,N_L)\in\mathbb Z^2\) and \(x,y\in\mathbb Z/L\mathbb Z\),
\begin{align}
 K_{\mathrm{CDW}}(x)^*K_{\mathrm{CDW}}(y)\delta_N
 &=\zeta^{(N_L-N_R-1)(y-x)}\delta_N,
 \label{eq:ch20b-klein-cdw}\\
 K_{\mathrm{SC}}(x)^*K_{\mathrm{SC}}(y)\delta_N
 &=\zeta^{(N_R+N_L)(y-x)}\delta_N.
 \label{eq:ch20b-klein-sc}
\end{align}
In particular \(K_J(x)^*K_J(x)=I\).
\end{lemma}
\begin{proof}
Write \(\chi_{\mathrm{CDW}}(N,x)
=\zeta^{(N_L-N_R-1)x}\) and
\(\chi_{\mathrm{SC}}(N,x)=\zeta^{(N_R+N_L)x}\).
The rightmost factor in \eqref{eq:ch20b-Kcdw} produces
\(\chi_{\mathrm{CDW}}(N,x)\).  Next, \(F_L\) lowers the left
charge and contributes \(P_L(N)\), after which \(F_R^*\) raises the
right charge and contributes \(P_R(N-e_L+e_R)\).  Thus
\[
 K_{\mathrm{CDW}}(x)\delta_N
 =\chi_{\mathrm{CDW}}(N,x)
   P_L(N)P_R(N-e_L+e_R)
   \delta_{N+e_R-e_L}.
\]
For the pairing word, both shifts lower charge, in the given order:
\[
 K_{\mathrm{SC}}(x)\delta_N
 =\chi_{\mathrm{SC}}(N,x)
   P_L(N)P_R(N-e_L)
   \delta_{N-e_R-e_L}.
\]
In each expression the product of the two signed Klein coefficients
is a real number of absolute value one, independent of \(x\).
The shift target is unique, so the adjoint at \(x\) takes that target
back to \(\delta_N\) and multiplies by the complex conjugate of its
coefficient at \(x\).  The two signs square to one, leaving
\(\overline{\chi_J(N,x)}\chi_J(N,y)\), which is precisely the
stated phase.  This computation explicitly preserves the source,
intermediate, and target charge labels.
\end{proof}

A formal exponential requires no analytic convergence.  Introduce a
\emph{new central indeterminate} \(t\) with \(t^*=t\).
The ring \(\mathfrak A[[t]]\) consists of all formal sequences
\(\sum_{n\geq0}a_nt^n\), \(a_n\in\mathfrak A\), with multiplication
by the finite convolution
\begin{equation}
 [t^n]\biggl(\sum_{j\geq0}a_jt^j\biggr)
        \biggl(\sum_{k\geq0}b_kt^k\biggr)
 =\sum_{j=0}^n a_jb_{n-j}.
 \label{eq:ch20b-convolution}
\end{equation}
There is no requirement that the coefficient algebra be commutative.
The notation \([t^n]\) means the coefficient of \(t^n\), not an
operator commutator.  For \(a\in\mathfrak A\) define
\begin{equation}
 \operatorname{Exp}_{\mathrm f}(ta)
 =\sum_{n=0}^\infty\frac{t^na^n}{n!}\in\mathfrak A[[t]],
 \qquad a^0=I.
 \label{eq:ch20b-formal-exp}
\end{equation}
Factorials are invertible \emph{complex scalar} coefficients.
More generally, if \(u\in t\mathfrak A[[t]]\) has zero constant
coefficient, then
\(\operatorname{Exp}_{\mathrm f}(u)=\sum_{k\geq0}u^k/k!\)
is a well-defined formal series: for each degree \(n\), only
\(k\leq n\) can contribute.  In particular,
\(\operatorname{Exp}_{\mathrm f}(\lambda t^2)\) for
\(\lambda\in\mathbb C\) denotes
\(\sum_{k\geq0}\lambda^kt^{2k}/k!\).
The positive functional \(\omega_{P,N}\) extends coefficientwise to a
\(\mathbb C[[t]]\)-linear evaluator
\begin{equation}
 \widehat\omega_{P,N}\biggl(\sum_{n\geq0}a_nt^n\biggr)
 =\sum_{n\geq0}\omega_{P,N}(a_n)t^n.
 \label{eq:ch20b-omega-hat}
\end{equation}
Although each original \(\omega_{P,N}\) is positive, this map is not
an ordinary scalar-valued state on arbitrary formal series, for which
there is no assumed coefficientwise nonnegative order.  Above all,
\(t\) is \emph{formal}: setting \(t=1\) is not defined by
\eqref{eq:ch20b-convolution}--\eqref{eq:ch20b-formal-exp}.

\begin{lemma}[Formal exponential product, with its actual hypothesis]
\label{lem:ch20b-exp-product}
If \(a,b\) belong to a unital complex algebra and \(ab=ba\), then
in its noncommutative-coefficient series ring
\begin{equation}
 \operatorname{Exp}_{\mathrm f}(ta)
 \operatorname{Exp}_{\mathrm f}(tb)
 =\operatorname{Exp}_{\mathrm f}(t(a+b)).
 \label{eq:ch20b-formal-binomial}
\end{equation}
In particular
\(\operatorname{Exp}_{\mathrm f}(-ta)
 \operatorname{Exp}_{\mathrm f}(ta)=I\).
No such factorization is claimed for noncommuting \(a,b\).
\end{lemma}
\begin{proof}
At degree \(n\), the product on the left has coefficient
\[
 \sum_{k=0}^n\frac{a^kb^{n-k}}{k!(n-k)!}
 =\frac1{n!}\sum_{k=0}^n\binom nk a^kb^{n-k}
 =\frac{(a+b)^n}{n!}.
\]
The last equality is the binomial formula, valid because \(ab=ba\).
Every coefficient is a finite sum, so no limit or rearrangement of
infinite analytic sums is involved.  Taking \(b=-a\) gives the inverse.
\end{proof}

The final ingredient is a Gaussian moment identity.  We prove it
rather than invoking a probability distribution on the infinite
polynomial space.
\begin{lemma}[Vacuum Wick recurrence and formal Gaussian identity]
\label{lem:ch20b-wick}
Let \(Y\) be a Hermitian \emph{linear combination} of dressed
creation and annihilation generators, with no constant term.  Put
\(V=\omega_P(Y^2)\).  Then \(V\geq0\),
\begin{align}
 \omega_P(Y^{2k+1})&=0,
 &\omega_P(Y^{2k})&=\frac{(2k)!}{2^k k!}\,V^k,
 \qquad k\geq0,
 \label{eq:ch20b-gauss-moments}\\
 \widehat\omega_P\bigl(\operatorname{Exp}_{\mathrm f}(itY)\bigr)
 &=\operatorname{Exp}_{\mathrm f}(-Vt^2/2)
 \quad\text{in }\mathbb C[[t]].
 \label{eq:ch20b-gauss-exp}
\end{align}
For \(k=0\), the even moment is understood to equal \(\omega_P(I)=1\).
\end{lemma}
\begin{proof}
Pull back by \(\beta_P^{-1}\) and denote the resulting operator
again by \(Y\).  Hermiticity allows the decomposition
\(Y=C_f+A_{\overline f}\), where
\(C_f=\sum_j f_jC_j\) and
\(A_{\overline f}=\sum_j\overline{f_j}A_j\), for finitely many
oscillator labels \(j=(\nu,m)\).  Consequently
\([A_{\overline f},Y]=V I\), where
\begin{equation}
 V=\sum_j m_j|f_j|^2
   =\omega_0(Y^2)\geq0.
 \label{eq:ch20b-GaussianV}
\end{equation}
The equality follows because only the ordered \(AC\) term survives
between vacuum vectors.  In the polynomial representation,
\(A_{\overline f}1=0\) and
\(\langle1,C_fq\rangle_0=0\) for all polynomials \(q\).
Since the commutator \([A_{\overline f},Y]\) is central, induction on
\(n\) gives \([A_{\overline f},Y^n]=nVY^{n-1}\).  Hence for
\(n\geq1\)
\begin{align*}
 \omega_0(Y^{n+1})
 &=\langle1,(C_f+A_{\overline f})Y^n1\rangle_0
 =\langle1,A_{\overline f}Y^n1\rangle_0\\
 &=\langle1,[A_{\overline f},Y^n]1\rangle_0
 =nV\,\omega_0(Y^{n-1}).
\end{align*}
Together with \(\omega_0(Y)=0\) and \(\omega_0(I)=1\), this
recurrence yields the claimed odd and even moments.  The pullback
\(\omega_P=\omega_0\circ\beta_P^{-1}\) restores the dressed formulation.
Finally, the coefficient of \(t^{2k}\) on the left of
\eqref{eq:ch20b-gauss-exp} is
\[
 \frac{i^{2k}}{(2k)!}\omega_P(Y^{2k})
 =\frac{(-V/2)^k}{k!},
\]
while all odd coefficients vanish.  These are exactly the coefficients
of the right-hand scalar formal exponential.
\end{proof}

\subsection{The formal vertices and their two-point functions}
\label{subsec:ch20b-correlators}

\begin{definition}[Weyl-type formal CDW and pairing vertices]
\label{def:ch20b-vertices}
For \(J\in\{\mathrm{CDW},\mathrm{SC}\}\), with \(K_J\) given in
\eqref{eq:ch20b-Kcdw}--\eqref{eq:ch20b-Ksc} and \(X_J\) given in
\eqref{eq:ch20b-Xcdw}--\eqref{eq:ch20b-Xsc}, define the formal series
\begin{align}
 W_{\mathrm{CDW}}(P,t,x)
 &=\frac1L\,F_R^*F_L\zeta^{(N_L-N_R-1)x}
    \operatorname{Exp}_{\mathrm f}(itX_{\mathrm{CDW}}(P,x)),
 \label{eq:ch20b-Wcdw}\\
 W_{\mathrm{SC}}(P,t,x)
 &=\frac1L\,F_RF_L\zeta^{(N_R+N_L)x}
    \operatorname{Exp}_{\mathrm f}(itX_{\mathrm{SC}}(P,x)).
 \label{eq:ch20b-Wsc}
\end{align}
These are \emph{Weyl-type formal vertices}, not normal-ordered
Mattis--Mandelstam fields.  The oscillator series commute with the
charge-word factors because they live in independent tensor factors.
\end{definition}

The choice not to normal-order is substantive.  Normal ordering of
creation and annihilation parts changes vacuum contractions and the
prefactor needed to identify a particular physical observable.  Neither
normal ordering nor exponentiation of \emph{compressed} field operators
can be inserted here without a new theorem.

\begin{theorem}[Coefficientwise formal vertex correlators; draft]
\label{thm:ch20b-correlators}
Assume \(L=2h\) with \(h\geq1\), \(M\geq1\) finite,
\(N\in\mathbb Z^2\), and real \(c,s\) satisfying \(c^2-s^2=1\).
On the constructed algebraic oscillator/charge carrier, the
coefficientwise evaluator satisfies
\begin{align}
 &\widehat\omega_{P,N}\!\left(
    W_{\mathrm{CDW}}(P,t,x)^*W_{\mathrm{CDW}}(P,t,y)\right)
 \nonumber\\[-2pt]
 &\hspace{1cm}=
 \frac{\zeta^{(N_L-N_R-1)(y-x)}}{L^2}
 \operatorname{Exp}_{\mathrm f}\!\left(-2gD_1(x,y)t^2\right),
 \label{eq:ch20b-cdw-corr}\\[4pt]
 &\widehat\omega_{P,N}\!\left(
    W_{\mathrm{SC}}(P,t,x)^*W_{\mathrm{SC}}(P,t,y)\right)
 \nonumber\\[-2pt]
 &\hspace{1cm}=
 \frac{\zeta^{(N_R+N_L)(y-x)}}{L^2}
 \operatorname{Exp}_{\mathrm f}\!\left(-2g^{-1}D_1(x,y)t^2\right).
 \label{eq:ch20b-sc-corr}
\end{align}
The exponentials on the right are scalar \emph{formal} series in
\(\mathbb C[[t]]\).  Equality means equality of every coefficient;
it does not involve an analytic value at \(t=1\).
\end{theorem}
\begin{proof}
We give each algebraic step, as the ordering is important.  The
complex conjugate of \(i\) is \(-i\), \(t^*=t\), and the \(X_J\)
are Hermitian.  Therefore
\begin{equation}
 \bigl(\operatorname{Exp}_{\mathrm f}(itX_J(P,x))\bigr)^*
 =\operatorname{Exp}_{\mathrm f}(-itX_J(P,x))
 \label{eq:ch20b-star-exp}
\end{equation}
coefficient by coefficient.  Adjunction reverses product order.  Since
every charge factor commutes with every oscillator factor,
\begin{align*}
 W_J(P,t,x)^*W_J(P,t,y)
 &=L^{-2}\operatorname{Exp}_{\mathrm f}(-itX_J(P,x))
      K_J(x)^*K_J(y)
      \operatorname{Exp}_{\mathrm f}(itX_J(P,y))\\
 &=L^{-2}K_J(x)^*K_J(y)
      \operatorname{Exp}_{\mathrm f}(-itX_J(P,x))
      \operatorname{Exp}_{\mathrm f}(itX_J(P,y)).
\end{align*}
Here we have \emph{not} commuted two arbitrary charge factors past
one another.  We moved only a tensor-factor charge product past an
oscillator product.

Lemma~\ref{lem:ch20b-hermitian-commuting} establishes
\([X_J(P,x),X_J(P,y)]=0\), so the
coefficientwise binomial identity
\eqref{eq:ch20b-formal-binomial} gives the \emph{exact} series identity
\begin{equation}
 W_J(P,t,x)^*W_J(P,t,y)
 =L^{-2}K_J(x)^*K_J(y)
   \operatorname{Exp}_{\mathrm f}(it\Delta X_J).
 \label{eq:ch20b-product-series}
\end{equation}
Apply \(\widehat\omega_{P,N}\) coefficient by coefficient.
Because the state is a tensor-product vector functional and the
charge operator is independent of the oscillator operators, every
coefficient factorizes into a charge expectation times an oscillator
expectation.  By Lemma~\ref{lem:ch20b-klein}, the charge factor equals
\(\zeta^{\alpha_J(N)(y-x)}\), with
\[
 \alpha_{\mathrm{CDW}}(N)=N_L-N_R-1,
 \qquad\alpha_{\mathrm{SC}}(N)=N_R+N_L.
\]
The Gaussian lemma applies to the Hermitian \emph{linear} oscillator
combination \(\Delta X_J\).  By
Proposition~\ref{prop:ch20b-variance}, its variance is respectively
\(V_{\mathrm{CDW}}=4gD_1(x,y)\) or
\(V_{\mathrm{SC}}=4g^{-1}D_1(x,y)\).  Thus the oscillator factor is
\(\operatorname{Exp}_{\mathrm f}(-V_Jt^2/2)\).  Substitution of the
two values gives \eqref{eq:ch20b-cdw-corr} and
\eqref{eq:ch20b-sc-corr}, including the factor \(L^{-2}\).
Every identity was established within a specified algebra and every
infinite-looking equality was checked coefficientwise.
\end{proof}

Two immediate consequences make the normalization explicit.
At \(x=y\), \(D_1(x,x)=0\) and the charge phase is one.  Thus
\begin{equation}
 \widehat\omega_{P,N}(W_J(P,t,x)^*W_J(P,t,x))=L^{-2}
 \quad\text{as a constant formal series.}
 \label{eq:ch20b-coincident}
\end{equation}
In particular, if one writes
\(\widehat\omega(W_J(x)^*W_J(y))=
  L^{-2}\chi_J(N;y-x)\sum_{k\geq0}C_k^Jt^{2k}\),
then \(C_0^{\mathrm{CDW}}=C_0^{\mathrm{SC}}=1\).  For any \(x,y\)
the remaining coefficients can be read directly:
\begin{equation}
 C_k^{\mathrm{CDW}}=\frac{(-2gD_1(x,y))^k}{k!},\qquad
 C_k^{\mathrm{SC}}=\frac{(-2g^{-1}D_1(x,y))^k}{k!}.
 \label{eq:ch20b-coefficients}
\end{equation}
They are complex-scalar coefficients of a formal series, not
measurements of a lattice decay exponent.

\subsection{Worked four-site example and a finite-CAR counterexample}
\label{subsec:ch20b-example}

Take \(L=4\), \(h=2\), \(M=1\), and
\(\zeta=i\), with positions \(x=0\), \(y=1\).
In the free Bogoliubov case \(c=1,s=0\), one has \(g=g^{-1}=1\).
For the source charge \(N=(0,0)\),
\[
 D_1(0,1)=1-\cos(\pi/2)=1,
 \quad\alpha_{\mathrm{CDW}}(0,0)=-1,
 \quad\alpha_{\mathrm{SC}}(0,0)=0.
\]
The theorem therefore specializes to
\begin{align}
 \widehat\omega_{P,(0,0)}(W_{\mathrm{CDW}}(t,0)^*
                                  W_{\mathrm{CDW}}(t,1))
 &=-\frac{i}{16}\operatorname{Exp}_{\mathrm f}(-2t^2),
 \label{eq:ch20b-example-cdw}\\
 \widehat\omega_{P,(0,0)}(W_{\mathrm{SC}}(t,0)^*
                                  W_{\mathrm{SC}}(t,1))
 &=\frac1{16}\operatorname{Exp}_{\mathrm f}(-2t^2).
 \label{eq:ch20b-example-sc}
\end{align}
To see the first nonconstant coefficient without trusting the closed
formula, compute \(V_J=4\) directly from the two chiral contributions.
The \(t^2\) coefficient of the oscillator expectation is
\[
 \frac{i^2}{2!}\,\omega_P((\Delta X_J)^2)
 =-\frac42=-2,
\]
as predicted.  Higher even coefficients are
\(2,-4/3,\ldots\), since
\begin{equation}
 \operatorname{Exp}_{\mathrm f}(-2t^2)
 =1-2t^2+2t^4-\frac43t^6+\cdots.
 \label{eq:ch20b-example-series}
\end{equation}
The only distinction between the two answers in this example is the
fixed, ordered Klein/charge phase.  The oscillator kernel is the same
because \(g=1\).

\begin{example}[Equal-point finite-CAR normalization is different]
\label{ex:ch20b-finite-counterexample}
On the \(L=4\) CAR Fock space of Section~20.1, choose a normalized
Slater sea \(\lvert\Omega\rangle\) with exactly two occupied momentum
orbitals in each of the independent \(R\) and \(L\) species.  The
normalized Fourier fields obey
\(\{c_{\nu,x},c_{\eta,y}^\dagger\}=
 \delta_{\nu\eta}\delta_{xy}\).
Translation invariance of a filled set of plane-wave orbitals gives
\[
 \langle\Omega\lvert n_{R,x}\rvert\Omega\rangle
 =\langle\Omega\lvert n_{L,x}\rvert\Omega\rangle=\frac24=\frac12,
 \qquad n_{\nu,x}=c_{\nu,x}^\dagger c_{\nu,x}.
\]
The even number operators on distinct species factorize in this
product Slater state.  Using the exact CAR words
\(O_{\mathrm{CDW}}=c_{R,x}^\dagger c_{L,x}\) and
\(O_{\mathrm{SC}}=c_{R,x}c_{L,x}\), we obtain
\begin{align*}
 \langle\Omega\lvert O_{\mathrm{CDW}}(x)^\dagger
               O_{\mathrm{CDW}}(x)\rvert\Omega\rangle
 &=\langle n_{L,x}(1-n_{R,x})\rangle_\Omega=\tfrac14,\\
 \langle\Omega\lvert O_{\mathrm{SC}}(x)^\dagger
               O_{\mathrm{SC}}(x)\rvert\Omega\rangle
 &=\langle n_{R,x}n_{L,x}\rangle_\Omega=\tfrac14.
\end{align*}
But \eqref{eq:ch20b-coincident} gives \(L^{-2}=1/16\) for either
formal Weyl-type vertex.  Therefore \emph{even at coincident points and
in this elementary free example}, these formally defined vertices are
not equal, as normalized observables and states, to the finite-CAR
bilinears.  The discrepancy is neither a numerical error nor a failure
of the theorem: the two sides were defined on different carriers and
with different operator normalization prescriptions.  The value
\(1/4\) is an example-specific finite-state calculation, not a
universal CAR identity.
\end{example}

\subsection{Vertex exchange, zero modes, and finite-cutoff approximations}
\label{subsec:ch20b-vertex-exchange}

We now examine how fermionic statistics are reconstructed from bosonic vertices.
Consider bare normal-ordered oscillator vertices $V_R(x) := :\!e^{i X_R(x)}\!:$ and
$V_L(y) := :\!e^{i X_L(y)}\!:$.
Because $[b_{R,n}, b_{L,m}^\dagger] = 0$, the two oscillator families commute.
By the normal-ordering multiplication law \eqref{eq:ch19-compression-remainder},
\[
 V_R(x) V_L(y) = :\!e^{i X_R(x) + i X_L(y)}\!: = V_L(y) V_R(x).
\]
Therefore, naive oscillator vertices are strictly \emph{commuting}:
\begin{equation}
 [V_R(x), V_L(y)] = 0, \qquad \{V_R(x), V_L(y)\} = 2 V_R(x) V_L(y) \ne 0.
 \label{eq:ch20b-naive-commuting}
\end{equation}
Bogoliubov dressing does not alter this: by \eqref{eq:ch19-cross-chiral-comm-zero},
$[A_R(x), A_L(y)] = 0$ as well.

\begin{theorem}[Mixed-Species Fermionic Anticommutation]
\label{thm:ch20b-mixed-car}
Let $\Psi_r(x) := F_r V_r(x)$ be the Klein-dressed vertex operators on
$\mathcal C \otimes \mathcal P$.  Then the mixed-species fields satisfy exact CAR:
\begin{equation}
 \{\Psi_R(x), \Psi_L(y)\} = 0, \qquad \{\Psi_R(x), \Psi_L^*(y)\} = 0.
 \label{eq:ch20b-mixed-car}
\end{equation}
\end{theorem}
\begin{proof}
Because oscillator modes and charge shifts live in independent tensor factors,
$F_R$ commutes with $V_L(y)$, and $F_L$ commutes with $V_R(x)$.  Using
$[V_R(x), V_L(y)] = 0$ and Theorem~\ref{thm:ch20b-shift-anticommute}:
\begin{align*}
 \Psi_R(x) \Psi_L(y) + \Psi_L(y) \Psi_R(x)
  &= F_R V_R(x) F_L V_L(y) + F_L V_L(y) F_R V_R(x) \\
  &= (F_R F_L + F_L F_R) V_R(x) V_L(y) = 0.
\end{align*}
The adjoint relation $\{\Psi_R(x), \Psi_L^*(y)\} = 0$ follows identically from $\{F_R, F_L^*\} = 0$.
\end{proof}

Next, consider \textbf{same-chirality exchange}.
For $x \ne y$, let $t = 2\pi(x-y)/L$ with $0 < t < 2\pi$.
Using the central commutator \eqref{eq:ch19-field-commutator}, Baker--Campbell--Hausdorff
multiplication yields
\begin{equation}
 V_R(x) V_R(y) = \exp\left( -2i \sum_{n=1}^M \frac{e^{-a q_n}}{n} \sin(q_n(x-y)) \right) V_R(y) V_R(x).
 \label{eq:ch20b-finite-osc-phase}
\end{equation}
In the continuum limit ($M \to \infty, a \to 0^+$), we use the Fourier sawtooth series:
\begin{equation}
 \sum_{n=1}^\infty \frac{\sin(n t)}{n} = \frac{\pi - t}{2} \qquad (0 < t < 2\pi).
 \label{eq:ch20b-sawtooth}
\end{equation}
The continuum oscillator exchange phase is therefore
\[
 \exp\left( -2i \frac{\pi - t}{2} \right) = e^{-i(\pi - t)} = -e^{it}.
\]
Notice that the oscillator phase $-e^{it}$ alone is \emph{not} equal to $-1$: it contains
an extra coordinate-dependent factor $e^{it}$.

The remaining factor is canceled by the zero-mode charge shift.
On the charge lattice $\mathcal C$, define the position-dependent zero-mode operator
\begin{equation}
 Z_R(x) := F_R \exp\left( i \frac{2\pi}{L}(N_R + \nu_R) x \right), \qquad \nu_R \in \{0, \tfrac{1}{2}\}.
 \label{eq:ch20b-zero-mode-field}
\end{equation}
Because $[N_R, F_R] = -F_R$, moving $F_R$ past $\exp(i \frac{2\pi}{L} N_R y)$ shifts $N_R \to N_R - 1$, yielding
\begin{equation}
 Z_R(x) Z_R(y) = e^{-it} Z_R(y) Z_R(x), \qquad t = \frac{2\pi(x-y)}{L}.
 \label{eq:ch20b-zero-mode-exchange}
\end{equation}
Combining the zero-mode and continuum oscillator factors for the full field $\Psi_R(x) = Z_R(x) V_R(x)$:
\begin{equation}
 \Psi_R(x) \Psi_R(y) = (-e^{it})(e^{-it}) \Psi_R(y) \Psi_R(x) = -\Psi_R(y) \Psi_R(x) \qquad (x \ne y \pmod L).
 \label{eq:ch20b-separated-fermion-exchange}
\end{equation}
The spin structure offset $\nu_R$ selects periodic ($\nu_R = 1/2$) or antiperiodic ($\nu_R = 0$) boundary conditions.

\begin{note}[Finite Mode Cutoff Warning]
\label{warn:ch20b-finite-cutoff-exchange}
Identity \eqref{eq:ch20b-separated-fermion-exchange} holds \textbf{only in the continuum limit}
$M \to \infty, a \to 0^+$.  At finite mode cutoff $M < \infty$, the finite Fourier sum
$\sum_{n=1}^M \frac{\sin(nt)}{n}$ differs from the sawtooth line $(\pi-t)/2$ by the Gibbs
oscillation residual:
\[
 \Delta_M(t) := \frac{\pi - t}{2} - \sum_{n=1}^M \frac{\sin(nt)}{n} \ne 0.
\]
Consequently, at finite $M$, the exchange of two right-moving vertex operators acquires
the non-fermionic phase factor $-e^{2i\Delta_M(t)} \ne -1$.
A finite-mode oscillator vertex is therefore \emph{an approximation to a fermion field},
not an exact fermionic operator.  Exact CAR on the finite lattice is provided exclusively
by the discrete Fourier CAR of Theorem~\ref{thm:ch20-discrete-fourier-car} and the
Jordan--Wigner occupation cocycle of Theorem~\ref{thm:ch20-constructive-car}.
\end{note}

\subsection{Scope and remaining comparison problem}
\label{subsec:ch20b-scope}

The results of this section are \emph{exact mathematical identities in
the constructed algebraic model}, subject to the displayed finite-mode,
real-parameter, charge-ordering, and state hypotheses.  Positivity of
\(\omega_{P,N}\) is a property of the uncompressed polynomial/charge
representation.  The series evaluator \(\widehat\omega_{P,N}\) is an
algebraic coefficientwise extension; it was never asserted to be a
convergent function of \(t\).  In particular there is no evaluation at
\(t=1\), no thermodynamic limit, and no power-law asymptotics in these
proofs.

Equally important, these formulas do not remove the projection
residuals of Section~20.1 or identify a compressed interacting finite-
CAR ground state with \(\omega_{P,N}\).  Such an identification would
require a precise operator dictionary, a comparison of the actual
states, control of cutoff and projection errors, and, if one wants an
analytic vertex, a topology and a convergence theorem.  The same
qualification applies to the duality construction proposed in
Section~\ref{sec:ch20-model-duality}: a star-algebra substitution may
transport this \emph{family} of formal states without being a unitary
symmetry of every finite-charge fermionic sector.

Finally, the relative normalization of raw charge/spin currents in a
later multicomponent theory cannot be changed silently by inserting a
\(1/\sqrt2\); and any use of the Fourier or Klein convention must
retain the stated orientation and source action.  Both conventions are
explicit here to make a future integration audit possible.  The formal-vertex
derivations in this section are included manuscript drafts awaiting independent
review and potential Lean formalization.  Section~20.1 has a separately accepted
finite-CAR theorem package in frozen Core; that acceptance does not extend to
the formal vertices or their comparison with an interacting finite-CAR state.
